\chapter{Traducción entre niveles}

\section{Traducción disciplinada}
\begin{definition}
Se llama traducción entre niveles al paso controlado mediante el cual una formulación, configuración o familia puede reexpresarse en otro nivel del canon sin perder su herencia estructural mínima.
\end{definition}

\begin{observation}
La traducción entre niveles no equivale a identidad plena entre niveles. Su función es preservar continuidad legible entre formulaciones no homogéneas.
\end{observation}

\section{Condición mínima de traducción}
\begin{definition}
Existe traducción disciplinada entre niveles cuando concurren simultáneamente:
\begin{enumerate}
  \item conservación suficiente de herencia canónica;
  \item compatibilidad de umbral para el propósito declarado;
  \item mantenimiento de una función reconocible;
  \item ausencia de contradicción estructural inmediata en el nivel de llegada.
\end{enumerate}
\end{definition}

\begin{observation}
La traducción no exige transportar toda la información disponible. Exige solo no perder aquella que resulte mínima para sostener el tratamiento disciplinado buscado.
\end{observation}

\section{Traducción y comparabilidad}
\begin{definition}
La traducción entre niveles puede utilizarse para hacer comparables piezas que de otro modo permanecerían separadas, siempre que la operación no borre sus diferencias operativas esenciales.
\end{definition}

\begin{observation}
Esta cláusula evita convertir la traducción en una reducción abusiva.
\end{observation}

\section{Traducción y familia}
\begin{definition}
Una familia operativa admite traducción entre niveles cuando sus configuraciones pueden conservar función común reconocible al pasar de un nivel a otro.
\end{definition}

\begin{observation}
No toda familia debe ser traducible de forma íntegra. Basta con que la traducción mantenga continuidad disciplinada suficiente para el alcance declarado.
\end{observation}

\section{Condición mínima de fallo}
\begin{definition}
Se considera fallida una traducción entre niveles cuando rompe herencia canónica, vuelve irreconocible la función de la pieza o introduce contradicción estructural en el nivel de llegada.
\end{definition}

\begin{observation}
Esta condición protege al canon intermedio frente a traslaciones solo retóricas o nominales.
\end{observation}

\section{Cláusula de alcance}
\begin{observation}
La traducción entre niveles fija una gramática mínima de paso entre formulaciones no homogéneas dentro del canon intermedio. No constituye todavía una teoría exhaustiva de equivalencias ni de representaciones finales del sistema.
\end{observation}

